% Tarea 1 - INF-221 Algoritmos y Complejidad - 2026-2
% Benjamin Lira - Rol: 202473586-1

Todas las mediciones se realizaron en el mismo equipo y con la misma
configuración:

\begin{center}
\begin{tabular}{ll}
\hline
Procesador        & Intel Core Ultra 7 155H (16 núcleos físicos, 22 lógicos) \\
Memoria           & 31,5 GB \\
Sistema operativo & Windows 11 Pro 10.0.26200 \\
Compilador        & \texttt{g++} (MSYS2 UCRT64) 15.2.0 \\
Opciones          & \texttt{-std=c++17 -O2 -Wall -Wextra -pedantic}, enlazado estático \\
Gráficos          & Python 3.14.3 con \texttt{numpy} 2.5.3, \texttt{matplotlib} 3.11.1 \\
\hline
\end{tabular}
\end{center}

Del cronómetro (\texttt{std::chrono::steady\_clock}) sólo se mide la llamada al
algoritmo: la lectura de los archivos, la copia del arreglo de trabajo y la
verificación de correctitud quedan fuera. La memoria reportada es el pico de
bytes vivos solicitados al \emph{heap}, medido interceptando \texttt{operator
new} y \texttt{operator delete}; no incluye la pila. Cada par (algoritmo, caso)
se ejecuta en un proceso propio, con una corrida de calentamiento no
cronometrada seguida de tres repeticiones. En los tamaños pequeños una sola
ejecución dura menos que unas pocas cuentas del reloj, por lo que el algoritmo
se ejecuta $K$ veces dentro de una única región cronometrada y se divide por
$K$, con las $K$ copias de trabajo preparadas de antemano.

Los conjuntos de prueba son los descritos en el enunciado. Para ordenamiento,
arreglos de $n \in \{10^1, 10^3, 10^5, 10^7\}$ elementos ---se añadió $10^6$
como punto intermedio para ajustar mejor la pendiente---, de tipo
\emph{ascendente}, \emph{descendente} o \emph{aleatorio}, con valores en
$D1 = \{0,\dots,9\}$ o $D7 = \{0,\dots,10^7\}$, y tres muestras por combinación.
Para multiplicación, pares de matrices de $n \in \{2^4, 2^6, 2^8, 2^{10}\}$, de
tipo \emph{dispersa}, \emph{diagonal} o \emph{densa}, con coeficientes en
$D0 = \{0,1\}$ o $D10 = \{0,\dots,9\}$, y tres muestras. Son 360 corridas de
ordenamiento y 144 de multiplicación. Todas se verificaron contra un oráculo
independiente del algoritmo medido, y todas resultaron correctas.
